\section{Research Proposal}

This section delineates the proposal to address the research problem identified in Section 5. 

This research proposes the development of a process-oriented framework to support the adoption, structural integration, and ongoing management of AI and LLM-based tools within Software Development processes. The proposal is motivated by the lack of structured and standardized guidance identified in the systematic literature review, particularly regarding how organizations can responsibly incorporate these technologies while preserving software quality and security, development integrity, and human oversight.

The proposed framework views AI and LLM-based tools as collaborative development services operating alongside human developers throughout the Software Development Lifecycle. Rather than positioning AI as an autonomous code generator, it emphasizes human–AI co-development, where AI systems provide computational and generative support while human professionals remain responsible for correctness, security, compliance, architectural soundness, and alignment with business requirements. Unlike traditional Agile and Scrum-based methodologies, which assume exclusively human contributors, the framework explicitly treats AI systems as first-class socio-technical actors whose contributions, limitations, and risks must be managed throughout the SDLC.

To operationalize this perspective, the framework defines a structured set of phases, practices, evaluation metrics, and decision-making guidelines governing the lifecycle of AI integration in Software Engineering:

\begin{enumerate}
    \item \textbf{Adoption \& Strategy Definition}, assessing justification, feasibility, and intended AI roles;
    \item \textbf{Design \& Integration}, establishing practices for human–AI collaboration, prompt formulation, model selection, and environment adaptation;
    \item \textbf{Operation \& Quality Assurance}, ensuring reliability, security, and transparency through human-in-the-loop mechanisms;
    \item \textbf{Continuous Monitoring \& Improvement}, applying metrics and governance to detect risks such as vulnerabilities, bias, or declining code diversity.
\end{enumerate}

Organizations increasingly seek methodological evolutions beyond traditional Agile and Scrum frameworks to address environments where AI actively contributes to software creation. The proposed framework represents such an evolution. While inspired by their iterative and value-driven principles, it is designed to replace human-only development methodologies in contexts where AI systems actively contribute to software creation, redefining process structure, responsibilities, and governance to explicitly account for human–AI collaboration. In practical terms, the framework aims to support AI-driven prototyping without sacrificing requirement accuracy, integrate AI responsibly into development and testing activities, redefine collaboration roles, reduce risks such as hallucinations and unverified code, and enable development processes in which human judgment and AI capabilities reinforce one another.

Thus, the proposal of this research is to define, structure, and validate a process-oriented framework that formalizes AI as a collaborative actor, introduces lifecycle governance mechanisms, and supports accountable human–AI collaboration in software development, providing actionable guidance for organizations seeking to adopt AI tools in a secure, responsible, sustainable and value-driven manner. This involves not only describing the framework but demonstrating it in practice and evaluating its suitability for real development contexts.
